\section{Data}\label{supp:data}

\subsection{Construction of the cell-week table}\label{supp:data-table}

The grid consists of all $0.5^\circ$ cells that intersect a country polygon of
Natural Earth (1:10m), excluding cells without a terrestrial biome (rock and ice, and
large lakes). Each cell is assigned to the country that contains its centroid. A
coastal cell whose centroid falls in the sea is assigned to the nearest country. Every
cell-week appears exactly once, and national indicators are attached through the
assigned country, with the exceptions described below. The result is 59{,}339 cells observed for up to 366 weeks.

The table built by the original pipeline differs in one important respect.
About 8.5\% of its 34.7 million distinct cell-weeks (before the biome exclusion)
occur more than once, with multiplicities such as $2^k$, $3^k$ and $2\cdot 3^k$, the
signature of many-to-many joins between the grid and polygons of countries, IPC
regions and national indicators. These copies make up 50.5\% of its 70.0 million
rows. The copies of a cell-week share its events but often disagree on the
covariates. Among duplicated cell-weeks, the copies differ in the EPI and the
liberal democracy index in 62\% of cases, in GDP in 49\%, in the IPC phases in about
25\% and in the climate variables in 7--11\%. We do not know why the climate values
of a cell differ between copies; the original pipeline does not record which
polygon each copy comes from. We keep the whole first row of each cell-week, so that
all covariates of a row come from the same copy. The climate values of duplicated
cell-weeks are therefore those of an arbitrary copy, and this part of the error is
retained rather than repaired. The climate covariates are controls in the joint
models and the conclusions do not rest on their coefficients. IPC
phase populations are the exception: they are averaged over the copies. Five national
indicators (EPI, liberal democracy index, greenhouse gas emissions, material and
ecological footprint) are re-attached for each country and year from the cells of
that country that were not duplicated. When a country-year has no value in these
cells, which happens when the indicator is missing for that country, the code falls
back to the most common value of the kept first copies of that country-year. Since
every country has non-duplicated cells, such values come from copies joined to a
neighbouring country's polygon. The fallback therefore can give a country its
neighbour's value and turns a missing value into an observed one, so the missingness
indicator is not set. In the whole repaired table, 14.5\% of rows take at least one
national indicator from this fallback. Access to electricity is not re-attached at
all: every cell-week keeps the value of its first copy, which in cells that had copies
can come from a neighbouring country. It is therefore not constant within a
country-year, and its decomposition into between- and within-country parts mixes differences across borders with change
over time. GDP per capita is a national quantity: World Bank GDP of the country
divided by the sum of the GHSL population of its cells, and it is constant within
each country-year. Section~\ref{sec:app-data} describes the consequences of the
duplication for the original analysis.

Two features of ACLED matter for modelling. First, geographic coverage expanded over
time, and several countries entered the dataset after 2018. Cell-weeks in a country
before its coverage start are \emph{unobserved}, not event free, so they are removed
rather than coded as zeros. Coverage start dates are taken from ACLED's country
coverage table. Territories not listed there take the date of the country or region
that contains them. Before the biome exclusion, coverage removes 7{,}078{,}822
cell-weeks in 79 countries and territories; 3{,}998{,}319 of them are in Antarctica,
which the biome exclusion would remove in any case. The event history is treated as
incomplete in the first four covered weeks. After both exclusions the table has
19{,}364{,}639 cell-weeks.
Second, the probability that an event is recorded depends on media and partner
density, so the outcome combines occurrence and reporting
\citep{weidmann2016closer}. The two are not separately identified
without auxiliary information. The latent fields of Section~\ref{sec:methods}
absorb spatially smooth reporting differences together with unmeasured drivers of
risk, and we interpret them accordingly. Reporting intensity also changes over time
and by country, for example when ACLED adds sources or partners. The year effects are
global, and at test time they are held at their 2022 value, so they cannot absorb
country-specific changes in reporting. Such changes may explain part of the rise in
recorded unrest above, and they can inflate both the estimated state dependence and
the persistence of the space-time fields.

ACLED also codes the spatial and temporal precision of each event. An event whose
exact location is unknown is placed at the centroid of the town or administrative
unit named in the source, and an event whose exact date is unknown at a nominal date
within the reported period. The data available to us are weekly counts by cell, and
they carry no precision codes, so we cannot filter or weight events by precision.
Imprecisely located events accumulate in the cells that contain administrative
centroids and capitals. This can inflate both the persistence of events in those
cells and the cell-level latent effects estimated there, and it is one more reason
to read those effects as recorded rather than actual risk.

\subsection{Covariates and missing values}\label{supp:data-cov}

The climate variables enter all models only through their means over the four weeks
before the current week, so that they precede the outcome; values of the current
week are not used. The original analysis also used the current week's values. A raw temperature
means different things in different climates. In the latent models, the five
climate variables (maximum and minimum temperature, precipitation, drought and
PM$_{2.5}$) therefore enter as standardised anomalies with respect to a cell- and
calendar-month-specific climatology estimated from the training years 2018--2022,
\[
z_{st} = \bigl(x_{st} - \bar{x}_{s,\mathrm{m}(t)}\bigr) / \mathrm{sd}_{s,\mathrm{m}(t)},
\]
where $\mathrm{m}(t)$ is the calendar month of week $t$. Precipitation is
transformed to $\log(1 + x)$ before the anomaly is computed. In dry seasons of arid
cells the training-period standard deviation can be zero or close to zero, and any
later rain would then give an extreme anomaly. We therefore floor
$\mathrm{sd}_{s,\mathrm{m}}$ at the 10th percentile of the positive cell-month
standard deviations of the same variable, and set $z_{st} = 0$ where the
climatological standard deviation is zero. In the AME region the floor applies, by
construction, to 10\% of the 164{,}628 cell-months for temperature and precipitation.
The drought index is constant over the training years in 138{,}600 cell-months (84\%),
so its anomaly is zero there, and the floor applies to a further 1.3\%. The drought
anomaly is therefore informative only in a minority of cell-months. The PM$_{2.5}$
anomaly is worse. Its climatological standard deviation is undefined in 91{,}392
cell-months (56\%), where PM$_{2.5}$ is missing, and the floor computed from the other
cell-months is about $10^{-9}$, so it does not prevent extreme anomalies. The
PM$_{2.5}$ anomaly is missing in 55.5\% of training rows, is then set to zero and has
no missingness indicator. We do not interpret its coefficient. The anomaly coefficients are identified
from variation over time within cells. They are therefore not confounded by the
time-invariant spatial pattern of the climatology. They can still be confounded by
regional shocks that vary over time, since the year effects are global. The cell's
mean over the training years enters separately as a cell-level covariate. A
five-year climatology is noisy, especially for precipitation, and a longer reference
period would be preferable where available. The anomalies are therefore best read
as departures from recent conditions.

\paragraph{Missing covariates.} Table~\ref{tab:covariates} hides substantial
missingness. In the repaired table over the AME region
(4{,}976{,}634 cell-weeks), the IPC phase populations are missing in 62.7\% of rows,
PM$_{2.5}$ in 55.5\%, and GDP in 3.6\%, mainly in South Sudan, Yemen,
Somaliland, Eritrea, Western Sahara and Syria. The other national indicators (EPI,
liberal democracy index, greenhouse gas emissions, material and ecological footprint)
are missing in 0.3\% of rows, at most about 850 cell-weeks in any country. These
shares are after the fallback described in Section~\ref{sec:data-acled}, which fills
some missing values with a neighbouring country's value. The
climate means over the previous four weeks are missing in 1.1\% of rows, and the event
histories in 1.1\%. Over all cell-weeks worldwide the IPC phases
are missing in 89.6\% of rows and precipitation in 20.8\%. IPC analyses are carried
out mainly where food crises are expected, and several of the countries without GDP
data are among the most conflict-affected in the region. Missingness is
therefore unlikely to be at random with respect to the outcome.

Apart from that fallback, we keep missing values as missing through the construction
of the features. In particular, a cell-week without an IPC analysis has a missing IPC value, not a
food-insecure population of zero. All covariates of the latent models other than the
event histories are standardised with training-period means and standard
deviations. Missing values are then set to zero, the training mean, and a
missingness indicator enters the model for each covariate group: one for each
national indicator and one for IPC (Section~\ref{sec:methods-model}).

\begin{table}
\caption{Covariates, resolution and source.}\label{tab:covariates}
\centering\small
\begin{tabular}{@{}lll@{}}
\toprule
Covariate & Resolution & Source \\
\midrule
Minimum / maximum temperature & cell--week & CHIRTS \\
Heat index, wet-bulb globe temperature & cell--week & CHIRTS \\
Precipitation & cell--week & CHIRPS \\
PM$_{2.5}$ & cell--week & CAMS \\
Drought index & cell--week & Global Drought Observatory \\
Events of each type in previous 4 weeks & cell--week & ACLED \\
Population & cell & GHSL \\
River length, lake area & cell & Natural Earth \\
Malaria risk & cell & Malaria Atlas Project \\
Biome & cell & RESOLVE Ecoregions \\
Food insecurity (IPC phases 1--5) & subnational--period & IPC \\
GDP per capita & country--year & World Bank \\
Access to electricity & country--year & World Bank \\
Environmental Performance Index & country--year & Yale EPI \\
Liberal democracy index & country--year & V-Dem \\
GHG emissions, material and ecological footprint p.c. & country--year & OWID, GFN \\
\bottomrule
\end{tabular}
\end{table}
